
\chapter{Introducción}
%% Este 'capitulo' o sección del documento no  numera ni secciones o subsecciones, se utiliza el esquema de libro para ordenar la información


\section{Definición del problema}

Uno de los problemas más comunes que se pueden encontrar en el país, tiene que ver con el manejo de los desechos sólidos. Según con el Diagnóstico Básico para la Gestión Integral de Residuos (DBGIR) realizado en el año 2020 por la Secretaría del Medio Ambiente y Recursos Naturales (SEMARNAT), se estima que una persona produce alrededor de 944 gramos de desechos al día. De esta forma, se calcula que al día se producen 120,128 toneladas de basura en el país. \cite{ANP2021} 

Dentro de este estudio, también se puede observar que los mayores generadores de residuos sólidos son el Estado de México, con 16,739 toneladas/día, y la Ciudad de México, con 9,552 toneladas/día, concentrando cerca del 21.88 \% de residuos nacionales por día. \cite{SEMARNAT2020}

En estos mismos datos se resalta como el 46.42 \% del total de desechos generados son orgánicos, otra cifra a destacar es que el 33 \% de los residuos orgánicos domésticos, son residuos alimenticios. \cite{SEMARNAT2020}

Sin embargo, las plantas compostadoras ubicadas en la Ciudad de México no son capaces de procesar toda la carga de residuos generados y la gran mayoría de esta materia orgánica es enviada a vertederos al aire libre, como los que se ubican en el Estado de México, lo que puede provocar daños al medio ambiente, así como un gasto en el transporte de la basura. \cite{SEMACDMX2021}

De esta forma, se plantea la fabricación de un sistema de compostaje que ayude a aprovechar los residuos orgánicos alimenticios. Esto porque son los residuos orgánicos que se generan con mayor frecuencia dentro de los hogares.

Para poder realizar este trabajo, ya se han realizado algunos proyectos previos que han venido ofreciendo la oportunidad de generar el compost de forma casera en poco tiempo; sin embargo, esta composta generada suele ser de mala calidad, debido a que deshidrata los residuos, perdiendo la capacidad de generar nutrientes suficientes para un jardín. 

Por otro lado, los compostadores que sí pueden generar el compost por medio de la descomposición suelen ser muy caros y debe de tomarse una capacitación especial para ser operados, lo que los vuelve un sistema no muy accesible para el público en general. 

De esta forma, se busca emplear procesos de ingeniería que permitan adaptar los procesos de compostaje para construir un dispositivo que genere compost y esté al alcance del público en general.

Para desarrollar este sistema, es necesario llevar a cabo un proceso de diseño que permita determinar las dimensiones y el aspecto necesarios para su operación, así como integrar dispositivos mecánicos, electrónicos, eléctricos e informáticos que trabajen en sinergia para controlar y monitorear cada una de las etapas del proceso de compostaje al interior del dispositivo. 

Por tal motivo, el problema específico que se pretende resolver es el diseño y construcción de un sistema semiautomático generador de compost para el aprovechamiento de los residuos orgánicos, el cual constará de una etapa de trituración, de agitación, monitoreo y el control de variables físicas. Como se ha mencionado, ese es el problema que se desea resolver; sin embargo, resulta evidente que durante el proceso de construcción de este proyecto se presentarán retos propios de la ingeniería.

Uno de los principales desafíos que se encuentran será el de poder controlar las variables físicas al interior del compostador, puesto que eso será determinante para la formación del compost. Por otra parte, una correcta elección de materiales y componentes influirá también en el control de dichas variables. Por último, no se puede dejar de lado el proceso de manufactura a seguir; elegir uno u otro influirá directamente en el costo que podría alcanzar el proyecto, así como el tiempo que puede tomar finalizarlo, buscando mantener en todo momento un diseño estético y cómodo para el usuario.
%%%%%%%%%%%%%%%

\section{Justificación}

En la actualidad, el manejo de los desechos urbanos es uno de los grandes retos a los que nos enfrentamos como sociedad. Ante esta problemática, han surgido diferentes estrategias que buscan aminorar nuestra huella ecológica y reducir el impacto de los residuos en el medio ambiente. Para contribuir a estos esfuerzos, proponemos fomentar la participación de la sociedad mediante el uso de un sistema mecatrónico. En este caso, se plantea la construcción de un sistema semiautomático para la generación de compost a partir de residuos orgánicos.

Al tratarse de una ingeniería multidisciplinaria la mecatrónica, incorpora elementos de distintas áreas, como la mecánica, control, informática y electrónica.

Para el presente proyecto, se buscará integrar las distintas ramas que la conforman, haciendo partícipe a la mecánica (diseño de elementos de máquinas y análisis estructural principalmente) en el diseño de la estructura. También se hará presente en el diseño del eje al interior del contenedor, que formará parte del sistema de agitación.

La electrónica se verá involucrada en el diseño y construcción de los circuitos que formarán parte de cada uno de los sistemas del compostador. También se aplicarán los conocimientos de la teoría de control para establecer las condiciones necesarias que permitan mantener las variables físicas del proceso de compostaje bajo ciertos parámetros que favorezcan su producción. Por último, serán necesarios conocimientos de informática, ya sea para la utilización de software de simulación y diseño, así como nociones de programación que permitan mejorar el funcionamiento del sistema.

Aunque de primera mano este proyecto parezca un tanto ambicioso, hasta cierto punto será innovador también, puesto que, como se podrá observar en el apartado de antecedentes, los proyectos existentes solo efectúan una parte de los procesos que el presente trabajo pretende realizar. Es decir, la mayoría de los proyectos encontrados solo realizan la trituración de la materia orgánica; otros se enfocan a mezclar los residuos. 

Muy pocos proyectos realizan el control de las variables físicas al interior del contenedor adicionando que la eliminación de los malos olores no es algo que la mayoría de los proyectos tomen en cuenta o resuelvan.

Escasos proyectos le dan importancia al tiempo de producción, ya sea del biogás o del abono, debido a que los procesos de descomposición que ocupan son diferentes. Para esta ocasión, se toma como referencia el método In-Vessel, que permite reducir el tiempo de producción que se puede tomar en cuenta.

Por último, es importante mencionar que el presente proyecto servirá como medio para fomentar el aprovechamiento de los residuos orgánicos, demostrando a las personas que es posible darle una segunda vida, un propósito a los alimentos una vez que estos se han convertido en basura.
%%%%%%%%%%%%%%%%%%%%

\section{Objetivo}

\subsection*{Objetivo General}

Diseñar e integrar un sistema semiautomático para la generación de compost a partir de residuos alimenticios, a través de un proceso de compostaje In-Vessel. \\

\addcontentsline{toc}{section}{Objetivos Particulares}  

\subsection*{Objetivos particulares de Trabajo Terminal I}
%5 − 7 objetivos
\begin{itemize}
    \item Diseñar la estructura del compostador para el procesamiento de compostaje In vessel.
    \item Diseñar en software CAD, los sistemas mecánicos para el proceso de compostaje.
    \item Simular con herramientas computacionales los módulos de control para el monitoreo de las variables físicas del proceso.
    \item Diseñar los circuitos electrónicos para los diferentes subsistemas del compostador.
    \item Desarrollar una interfaz humano-máquina que permita una comunicación con el usuario.
    \item Validar por medio de simulación con herramientas computacionales de cada uno de los módulos del sistema. 
    \item Integrar los módulos y verificar su funcionamiento con ayuda de herramientas computacionales.

\end{itemize}

\subsection*{Objetivos particulares de Trabajo Terminal II}

%3 − 7 objetivos
\begin{itemize}
    \item Fabricar los elementos mecánicos del sistema de compostaje por medio de diversos procesos de manufactura.
    \item Implementar el módulo de control con ayuda de una tarjeta de desarrollo y los sensores necesarios.
    \item Implementar la interfaz humano-máquina a través de la integración de esta con los componentes electrónicos para la operatividad del sistema.
    \item Verificar el correcto funcionamiento de cada uno de los módulos que conforman al compostador.
    \item Integrar y verificar el sistema de compostaje por medio de pruebas para la obtención de compost al final del proceso. 

\end{itemize}
%%%%%%%%%%%%%%%%%

\section{Enfoque mecatrónico}



%%%%%%%%%%%%%%%%%%%%%
\section{Antecedentes}
En el marco del desarrollo del proyecto, se han revisado diversos trabajos dentro de esta área de estudio (Tabla \ref{tab:Tab_Antecedentes}), a partir de los cuales se pueden referenciar diversos aspectos  de los subsistemas diseñados.

\begin{table}
    \centering
\caption{Antecedentes}
\label{tab:Tab_Antecedentes}
\resizebox{!}{19cm}{
\begin{turn}{90}
     \begin{tabular}{|m{3cm}|m{2cm}|m{5cm}|m{3cm}|m{3cm}|m{3cm}|} \hline 
         Proyecto&  Autor&  Descripción&  Aportes&  Ventajas& Desventajas\\ \hline 
         Estudio de volteadoras para mejorar la produccion de compost en la floricultura Nevado Ecuador de la Parroquia Mulalillo&  Eddisson Antonio Lozada Vinces&  Mecanismo que permite realizar un voltero de las pilas de compost, donde se controlan algunos parametros como la oxigenación, temperatura, humedad, olor y tiempo del proceso&  Bases de los sistemas de monitoreo y principio del proceso de aireación &  Caracteristicas óptimas de transporte, sistema de ensamblaje intuitivo y amigable con el usuario. & Capacidad de volteo de 220-825 $m^3 / h$ y solo realiza el volteo \\ \hline 
         Lomi&  &  Dispositivo que transforma los residuos de comida en compost, utilizando abrasión, calor y oxigeno&  Es de tamaño compacto, amigable con el usuario y tiene diversos tipos de operación&  Tamaño compacto, cuenta con 3 modelos de operacion (Expres, Eco y bioplasticos). & Capacidad de procesar hasta 3 L de materia. Tiene un costo muy elevado\\ \hline 
         Oklin&  &  Es un dispositivo que mediante el uso de tecnologia microbiana, reduce en un 90 \% el volumen de residuos alimenticios en compost, en 24 horas.&  Se implemte el uso de los mocroornaismos, llamados acidulos para acelerar el proceso. Se controla el ambiente para eliminar a bacterias dañinas.&  Es amigable con el usuario, tiene un tiempo de procesamiento corto. Cuanta con un sistema desodorante y cuenta con sistemas de higienización& Require un tiemo de espera la dapatacion de los microorganismos, tambien una manguera de drenaje y tiene un costo elevado\\ \hline 
    \end{tabular}   
\end{turn}
}
\end{table}

\begin{table}
    \centering
\caption{Antecedentes pt. 2}
\label{tab:Tab_Antecedentes_2}
\resizebox{!}{16cm}{
\begin{turn}{90}
     \begin{tabular}{|m{3cm}|m{2cm}|m{5cm}|m{3cm}|m{3cm}|m{3cm}|} \hline 
         Proyecto&  Autor&  Descripción&  Aportes&  Ventajas& Desventajas\\ \hline 
         Zera&  &  Como medio para acelerar el proceso de descomposición, este dispositivo combina el manejo de oxigeno, humedad y temperatura&  Conexión a red WIFI y uso de aplicación para control remoto&  Reduce el desperdicio a dos tercios de su volumen. Contenedor del tamaño de un bote de basura. Conexión WIFI y aplicación amigable con el usuario.& Producto descontinuado y no cumplia con lo prometido\\ \hline 
         Sistema de composta controlado por un sistema electrónico&  David Giovanni Rodríguez Chávez&  Contenedor capaz de regular la temperatura y humedad al interior del mismo, mediante el uso de un Arduino, sensor de temperatura, dos calentadores y una pantalla LCD para monitoreo de las mediciones.&  Es posible usar plastico en el contenedor, sin embargo existe riesgo de falla por el calor al interior&  Diseño amigable y economico de cada uno de sus subsistemas& Los materiales con los que está construido, pueden fallar. Escapan los malos olores. El uso del Arduino, no puede optimizar el proceso de lectura de los sensores.\\ \hline 
         Prototipo automático de alimentación de materia prima para un biodigestor.&  Jonathan Sánchez Maravilla&  Sistema de trituración de materia orgánica, que mezcla con agua la materia ingresada para mantener un valor de humedad constante y posteriormente almacenar la mezcla.&  Dispositivo eficiente y capaz de triturar huesos de res.&  Puede triturar materiales bastantes duros (huesos de res y cerdo) y mantiene un valor constante de sólidos al interior del contenedor& Requiere grandes cantidades de agua, solo representa la etapa de trituración\\ \hline 
    \end{tabular}   
\end{turn}
}
\end{table}

\begin{table}
    \centering
\caption{Antecedentes pt. 3}
\label{tab:Tab_Antecedentes_3}
\resizebox{!}{18cm}{
\begin{turn}{90}
     \begin{tabular}{|m{3cm}|m{2cm}|m{5cm}|m{3cm}|m{3cm}|m{3cm}|} \hline 
         Proyecto&  Autor&  Descripción&  Aportes&  Ventajas& Desventajas\\ \hline 
         Sistema biodigestor anaeróbico para hogares alimentado con heces caninas y residuos orgánicos de cocina&  Juan Sebastián González Diana&  Biodigestor de bajo costo que permite producir biogás y abono a partir de residuos orgánicos y heces de animales.&  Producción de boigás a un bajo costo y proceso de heces caninas&  Bajo costo, fácil implementación. Mide la temperatura, presión, concentración de gás y pH. & Requiere de una computadora protatil para el monitoreo de las mediciones, no cuenta con etapa de trituración, requiere fermentación adicional para generar humus, no realiza el control de las variables físicas y no resuelve el problema de los malos olores.\\ \hline 
         Prototipo automático de alimentación de materia para un biodigestor&  Giovanni Sinhue Camacho Ruiz&  Sistema que permite triturar materia orgánica, así mezclarla con agua en un dispositivo de almacenamiento manteniendo un valor de sólidos constantes.&  Sistema de medicion del nivel de la mezcla.&  Accesible económicamente, fácil de usar y es contruido con materiales accesibles.& Requiere de contidades considerables de agua para su funcionamiento. No realiza el control de las variables ni maneja tiempos de fermentación.\\ \hline  
    \end{tabular}   
\end{turn}
}
\end{table}

\begin{table}
    \centering
\caption{Antecedentes pt. 4}
\label{tab:Tab_Antecedentes_4}
\resizebox{!}{18cm}{
\begin{turn}{90}
     \begin{tabular}{|m{3cm}|m{2cm}|m{5cm}|m{3cm}|m{3cm}|m{3cm}|} \hline 
         Proyecto&  Autor&  Descripción&  Aportes&  Ventajas& Desventajas\\ \hline 
         Diseño de una máquina domestica automática para generar compost a partir de residuos orgánicos&  Renzo Rogger Acosta Gonzales&  Máquina doméstica capaz de generar compost, donde, por medio de una cámara multipropósitos se pueden realizar distintas tareas, así como el control de los párametros de compostaje, con un tiempo de producción de 5 dias.&  Cuenta con una cámara multipropósitos (Contiene, transporta y separa los residuos orgánicos).&  Monitoreo del proceso, control de variables físicas (temperatura y concentración de oxigeno), establece un tiempo de producción de 5 dias, permite ingresar desperdicios diariamente sin afectar el proceso de compostaje.& No cuenta con etapa de trituración, precio elevado, no resuelve el problema de los malos olores, oportunidad de mejora en la interfaz de usuario y celda de carga. \\ \hline 
         Caracterización del proceso de compostaje y aprovechamiento del calor generado en un reactor bajo aireación forzada.&  Diego Fallas Conejo.&  Sistema de compostaje de pulpa de café tipo reactor con control de aireación, así como los factores que rigen el proceso, ademas de establecer un aislamiento físico entre la materia prima y el ambiente.&  Utilización de un ventilador para aireación del sistema y estudio del comportamientode la materia con respecto a la aireación.&  Fácil de construir y uso de distintas tasas de aireación& La pulpa de café, al tener un alto contenido de humedad, presenta limitantes en el proceso (oxigenación principalmente) y oportunidad de mejora en el proceso de reducción de humedad.\\ \hline 
    \end{tabular}   
\end{turn}
}
\end{table}
%%%%%%%%%%%%%%%%%%%%%%%%%%%
\section{Capitulado}    %%  descripción general del documento mencionando contenidos de cada capitulo




%%%%%%fin del archivo
\endinput 